\begin{abcliste}
\abc The Compton wavelength is
\begin{align}
 \lambda_C &= \lCF = \frac{\nch}{\ncme\times\ncc} = \lC \approx \lCP{-12}{3}
\end{align}
and so
\begin{align}
 \theta = \angaO:&\quad \Delta\lambda = \lCP{-12}{3}\times\left(1-\cos\angaO\right) \approx \result{\dlaP{-12}{2}} \\
 \theta = \angbO:&\quad \Delta\lambda \approx \result{\dlbP{-12}{2}} \\
 \theta = \angcO:&\quad \Delta\lambda \approx \result{\dlcP{-12}{2}} \\
 \theta = \angdO:&\quad \Delta\lambda \approx \result{\dldP{-12}{2}}
\end{align}
$1-\cos\theta$ grows with the angle all the way to \SI{180}{\degree} (it does not stop at \SI{90}{\degree}): the change is largest at $\result{\angdO}$. A photon scattered straight back would have the largest change of all, $2\lambda_C$.

\abc The wavelength itself does not appear in the formula: $\Delta\lambda = \dlcP{-12}{2}$ is $\result{\text{the same for both}}$.

\abc The energy of a photon is $E = \frac{h c}{\lambda}$, so the scattered photon keeps the fraction $\frac{E'}{E} = \frac{\lambda}{\lambda+\Delta\lambda}$ and loses the fraction $\frac{\Delta\lambda}{\lambda+\Delta\lambda}$:
\begin{align}
 \text{X-rays:}&\quad \fxF = \frac{\dlcP{-12}{3}}{\lxO+\dlcP{-12}{3}} = \fx \approx \fxpP{0}{2} \\
 \text{gamma rays:}&\quad \fgF = \frac{\dlcP{-12}{3}}{\lgmO+\dlcP{-12}{3}} = \fg \approx \fgpP{0}{2}
\end{align}
The same $\Delta\lambda$ is a much larger part of a short wavelength: the $\result{\text{gamma photons}}$ lose the larger fraction of their energy. For visible light (\SI{500}{\nm}) the shift is far too small to notice: that is why the Compton effect needs X-rays or gamma rays.
\end{abcliste}
